\documentclass{article}

% --- Track option: preprint (arXiv-style, author names shown) ---
\usepackage[preprint]{neurips_2026}

% --- Required packages ---
\usepackage[utf8]{inputenc}
\usepackage[T1]{fontenc}
\usepackage{hyperref}
\usepackage{url}
\usepackage{booktabs}
\usepackage{amsfonts}
\usepackage{amsmath}
\usepackage{amssymb}
\usepackage{nicefrac}
\usepackage{microtype}
\usepackage{xcolor}
\usepackage{graphicx}

% --- Optional packages ---
\usepackage{float}

% --- Title & authors ---
\title{Size-Scaled AlphaZero Self-Play for Hex:\\A Single-Network Agent Across Board Sizes 5 to 17}

\author{
  Chorn Philipp \And
  Hromada Stefan \And
  Wang Li Wen \And
  Cahya Wirawan \\
  University of Applied Sciences Technikum Wien (FHTW) \\
  \texttt{\{ai25m0xx, ai25m0xx, ai25m0xx, ai25m063\}@technikum-wien.at}
}

\begin{document}

\maketitle

\begin{abstract}
Mastering the board game Hex with self-play reinforcement learning is challenging
because the branching factor and the length of the shortest winning connection both
grow quadratically with board side length, so a recipe tuned for a small board
collapses into near-random play on a larger one. We present a self-contained
AlphaZero-style agent for Hex that addresses this through a single design principle:
every search and learning quantity is derived automatically from the board size.
The simulation budget, residual-network width and depth, self-play volume, replay
capacity, and mini-batch size are all scaled by closed-form rules so that Monte Carlo
Tree Search visits each root child enough times for its visit-count distribution to
carry a usable policy target. A colour-agnostic perspective frame lets one network
play both players, and a compiled C++ tree-search kernel with batched, parallel leaf
evaluation keeps the GPU saturated. We further integrate five engineering refinements
common to strong implementations: rhombus symmetry augmentation, sub-tree reuse across
moves, calibrated resignation, mixed-precision training, and decoupled weight decay
with a cosine schedule. Trained from scratch on a single GPU, the agent reaches a
$90$--$100\%$ win rate against a uniform-random opponent on $11\times11$ and $13\times13$
boards within roughly $40$ self-play iterations, and continues to improve on
$15\times15$. We describe the scaling rules, the search engine, and the training loop in
full, and discuss what is needed to move from defeating random play to matching
specialist Hex engines.
\end{abstract}

\section{Introduction}

Hex is a two-player connection game invented independently by Piet Hein and John Nash.
Players alternately place stones on the cells of an $n\times n$ rhombus; one player
wins by forming an unbroken chain of their colour between their two opposing edges.
The game is elegant for theory and demanding for practice: it can never end in a draw,
the first player provably has a winning strategy by a strategy-stealing argument
\citep{nash1952hex}, yet that strategy is unknown for all but the smallest boards.
Hex has therefore long served as a benchmark for game-playing artificial intelligence
\citep{anthony2017thinking,gao2017move}, and it is the basis of the reinforcement
learning assignment at FHTW for which this work was developed: students implement an
agent behind a fixed engine interface and play it against the provided environment.

The dominant paradigm for games of this kind is AlphaZero \citep{silver2018general},
which couples Monte Carlo Tree Search (MCTS) with a deep neural network trained purely
on self-play. The network supplies move priors and position values to guide the search,
while the visit counts the search produces become a stronger policy target than the raw
network output. This mutual reinforcement---search improving the targets, the network
improving the search---is what drives learning without any human games or heuristics.
The recipe is conceptually simple but notoriously sensitive: too few simulations and
the visit-count policy target is dominated by noise, the bootstrap never ignites, and
the agent plateaus at random play.

This sensitivity is sharpened on Hex because difficulty scales steeply with the board.
The root branching factor is $n^2$, and the minimum number of stones needed to span the
board grows linearly in $n$, so both the breadth and the depth of meaningful search grow
with size. A fixed simulation count that is generous on a $5\times5$ board spreads only a
fraction of a visit across each child on a $15\times15$ board, where the root has $225$
legal moves. A practitioner who hard-codes hyperparameters tuned for a small board and
simply increases $n$ will watch the agent fail to learn, with no error to diagnose.

We address this with a deliberately small idea applied consistently: \emph{let the board
size determine everything}. From a single integer $n$ we derive the per-move simulation
budget, the network's channel count and residual depth, the number of self-play games per
iteration, the replay-buffer capacity, and the mini-batch size, each from a closed-form
rule motivated by how that quantity must grow to keep learning healthy. Around this core
we build a complete, reproducible system: a colour-agnostic perspective frame so one
network plays both sides, a compiled C++ search kernel with virtual-loss parallelism and
batched leaf evaluation, and five refinements drawn from the modern AlphaZero literature.

In summary, our contributions are:
\begin{itemize}
  \item We derive and justify a set of \emph{board-size scaling rules} that jointly set
        the MCTS budget, network capacity, and data volume, and we explain why each must
        grow with $n$ for the self-improvement loop to bootstrap.
  \item We describe a complete, single-GPU AlphaZero implementation for Hex featuring a
        perspective-frame single network, a compiled C++ MCTS engine with parallel
        virtual-loss leaf collection, and pipelined multi-process self-play.
  \item We integrate and document five practical refinements---symmetry augmentation,
        sub-tree reuse, calibrated resignation, mixed precision, and AdamW with cosine
        annealing---and report empirical learning curves across board sizes
        $11\times11$ through $15\times15$.
\end{itemize}
The remainder of the paper reviews related work (Section~\ref{sec:related}), details the
method and scaling rules (Section~\ref{sec:method}), describes the experimental setup
(Section~\ref{sec:experiments}), reports results (Section~\ref{sec:results}), interprets
them (Section~\ref{sec:discussion}), and concludes (Section~\ref{sec:conclusion}).

\section{Related Work}
\label{sec:related}

\subsection{Self-play search-and-learn agents}
AlphaGo first combined supervised policy learning, MCTS, and reinforcement learning to
defeat a professional Go player \citep{silver2016mastering}. AlphaGo Zero removed human
data, learning entirely from self-play with a single residual network producing both
policy and value \citep{silver2017mastering}, and AlphaZero generalised the method to
chess and shogi \citep{silver2018general}. MuZero further removed the need for a known
environment model by learning a latent dynamics model jointly with the value and policy
\citep{schrittwieser2020mastering}. Our work is a faithful instance of the AlphaZero
template rather than a new algorithm; our concern is the engineering and hyperparameter
discipline needed to make that template learn reliably on Hex across a range of board
sizes on modest hardware.

\subsection{Reinforcement learning for Hex}
Hex has a rich history as an AI testbed. Expert Iteration (ExIt) was introduced and
evaluated on Hex, framing the search-and-learn loop as an imitation of a slow expert by a
fast apprentice \citep{anthony2017thinking}. Deep convolutional networks were shown to
predict expert Hex moves and to strengthen the strong solver-based engine MoHex
\citep{gao2017move}, whose lineage rests on the connection-based theory and inferior-cell
analysis catalogued by \citet{hayward2019hex}. These specialist systems combine learned
evaluation with Hex-specific knowledge such as virtual connections and dead-cell pruning.
Our agent deliberately uses \emph{no} Hex-specific knowledge beyond the win condition and
a board symmetry; it learns connectivity implicitly through the value head. This makes it
weaker than knowledge-rich engines but cleanly isolates what a generic AlphaZero pipeline
can achieve, and clarifies the gap that domain knowledge would later fill.

\subsection{Monte Carlo Tree Search and its parallelisation}
MCTS with the UCT selection rule \citep{kocsis2006bandit} and the broader family surveyed
by \citet{browne2012survey} underpins modern game agents. AlphaZero replaces UCT's
exploration term with PUCT, in which a learned prior biases exploration toward promising
moves. Parallelising MCTS to feed a GPU efficiently requires care: virtual loss
\citep{chaslot2008parallel} temporarily penalises in-flight nodes so concurrent
simulations diverge to different branches, enabling many leaves to be evaluated in one
batched network call. Large-scale systems such as ELF OpenGo \citep{tian2019elf} and
KataGo \citep{wu2019accelerating} demonstrate how batching, sub-tree reuse, and
playout-cap or resignation tricks raise self-play throughput by orders of magnitude.
KataGo in particular contributes efficiency techniques---auxiliary targets, playout caps,
and aggressive batching---several of which inform the refinements we adopt.

\subsection{Value-based and policy-gradient baselines}
Before self-play search, the standard reinforcement learning toolkit comprises tabular
Q-learning \citep{watkins1992qlearning}, Deep Q-Networks \citep{mnih2015human},
the REINFORCE policy gradient \citep{williams1992simple}, and Proximal Policy
Optimization \citep{schulman2017proximal}. These methods learn a direct mapping from
state to action or value without lookahead at decision time. On a game with Hex's
combinatorial depth they struggle to assign credit through long connection sequences,
which is precisely the weakness that decision-time search remedies. Our repository
includes these methods as baselines, and their comparative weakness is the empirical
motivation for the search-based agent that is the focus of this paper.

\paragraph{Research gap.} The AlphaZero recipe is well documented for fixed, large
problems with industrial compute, and Hex-specific engines lean heavily on domain
knowledge. What is comparatively under-documented is how to make a \emph{knowledge-free}
AlphaZero agent learn \emph{robustly across a range of board sizes} on a single GPU,
where the same code must succeed on a $7\times7$ board and a $17\times17$ board without
manual retuning. This paper addresses that gap with explicit, justified size-scaling rules
and a complete reproducible implementation.

\section{Method}
\label{sec:method}

\subsection{Game representation and the perspective frame}
The engine represents the board as an $n\times n$ integer grid with values
$\{+1,-1,0\}$ for the two players and empty. Red connects left to right, Blue connects
top to bottom. Rather than train separate networks or feed the side-to-move as an extra
plane, we adopt a \emph{perspective frame}: when it is Blue's turn, the board is recoded
by transposing along the short diagonal and swapping colours, so that from the network's
point of view the player to move is always ``Red, connecting left to right.'' A chosen
move is mapped back to true coordinates by the same self-inverse transform. Every stored
state is thus in a canonical frame, which both halves the function the network must learn
and makes data augmentation straightforward.

The network input is a $2\times n\times n$ tensor whose two channels hold, respectively,
the stones of the player to move and the stones of the opponent in the perspective frame.

\subsection{Network architecture}
The agent uses a ResNet-style actor--critic, \texttt{HexAlphaNet}. A $3\times3$
convolutional stem with batch normalisation lifts the two input planes to $C$ channels,
followed by $B$ residual blocks, each two $3\times3$ convolutions with batch normalisation
and a skip connection \citep{he2016deep,ioffe2015batch}. Two heads share this trunk: a
policy head ($1\times1$ convolution to two planes, then a linear map to $n^2$ move logits)
and a value head ($1\times1$ convolution to one plane, a $256$-unit hidden layer, and a
$\tanh$ output in $[-1,+1]$). Invalid moves are masked by adding $-\infty$ to their logits
before the softmax.

\subsection{Search: PUCT with batched parallel evaluation}
At each move the agent runs $N$ MCTS simulations. Each simulation selects a path from the
root by maximising the PUCT score
\begin{equation}
  \mathrm{PUCT}(s,a) \;=\; -Q(s,a) \;+\; c_{\text{puct}}\, P(s,a)\,
  \frac{\sqrt{\sum_b N(s,b)}}{1+N(s,a)},
  \label{eq:puct}
\end{equation}
where $N(s,a)$ is the child visit count, $P(s,a)$ the network prior, and $Q(s,a)$ the mean
value stored \emph{from the child's player's perspective}; the leading negation converts it
to the selector's perspective, so we prefer children where the opponent fares worst. A
leaf is expanded by one batched network evaluation that supplies priors for its children
and a scalar value $v$ for the leaf. Values are propagated to the root under the negamax
convention, negating at each level:
\begin{equation}
  W(s,a) \leftarrow W(s,a) + v, \qquad
  Q(s,a) = \frac{W(s,a)}{N(s,a)}, \qquad v \leftarrow -v .
  \label{eq:backup}
\end{equation}
Because a terminal Hex node always means the player to move has just lost, terminal leaves
receive $v=-1$, which propagates correctly under Eq.~\eqref{eq:backup}.

To use a GPU efficiently we collect \emph{several} leaves before each network call. During
selection a \emph{virtual loss} is added to every node on the in-flight path---raising its
$Q$, lowering its $-Q$, and thus its PUCT score---so that subsequent concurrent selections
diverge to different branches; the virtual loss is undone during backup
\citep{chaslot2008parallel}. The number of leaves gathered per batch, $p$, must satisfy
$p \le N/7$ so that at least seven sequential search rounds occur per move; otherwise the
tree never grows beyond a couple of levels and the visit-count policy is nearly uniform.
At the root we mix Dirichlet noise into the priors,
$P \leftarrow (1-\varepsilon)P + \varepsilon\,\eta$ with
$\eta\sim\mathrm{Dir}(\alpha)$, to ensure exploration. Self-play games batch this search
across all concurrent games and, optionally, across multiple CPU worker processes, so a
single network forward pass serves many positions at once.

A compiled C++ extension (\texttt{hex\_mcts}) implements tree storage and traversal,
exposing a batched \texttt{select\_leaves}/\texttt{expand\_and\_backup} interface to
Python; the pure-Python search is retained as a drop-in fallback. The C++ kernel removes
per-node Python overhead, which dominates run time once the tree holds thousands of nodes.

\subsection{Training objective}
Self-play produces tuples $(s,\pi,z)$ where $\pi$ is the normalised root visit-count
distribution and $z\in\{+1,-1\}$ records whether the player to move at $s$ eventually won.
The network is trained to match both targets:
\begin{equation}
  \mathcal{L}(\theta) \;=\;
  \underbrace{-\textstyle\sum_a \pi(a)\,\log p_\theta(a\mid s)}_{\text{policy cross-entropy}}
  \;+\;
  \underbrace{\big(v_\theta(s) - z\big)^2}_{\text{value MSE}}
  \;+\; \lambda\,\lVert\theta\rVert_2^2 ,
  \label{eq:loss}
\end{equation}
with the $L_2$ term realised through decoupled weight decay in AdamW
\citep{loshchilov2019decoupled} rather than added to the loss explicitly. The learning
rate follows a cosine schedule over the planned iterations. The boxed procedure in
Figure~\ref{alg:az} summarises one training iteration.

\begin{figure}[H]
  \centering
  \fbox{\begin{minipage}{0.93\linewidth}
  \textbf{Algorithm 1: One AlphaZero training iteration for Hex.}\\[2pt]
  \textbf{Input:} network $\theta$, board size $n$, sims $N$, batch leaves $p$.
  \begin{enumerate}\setlength{\itemsep}{1pt}
    \item For each of $G(n)$ self-play games (batched together), until terminal:
      \begin{enumerate}\setlength{\itemsep}{1pt}
        \item Run $N$ PUCT simulations (Eq.~\ref{eq:puct}), collecting $p$ leaves per batched evaluation.
        \item Set $\pi \leftarrow$ normalised root visit counts; store $(s,\pi,\text{player})$.
        \item Sample a move from $\pi^{1/\tau}$ ($\tau$-greedy after move $T_{\text{cut}}$) and advance.
        \item Reuse the chosen child's sub-tree as the new root.
      \end{enumerate}
    \item Label every stored $s$ with $z=\pm1$ from the outcome; augment by $180^\circ$ symmetry.
    \item Push samples to the replay buffer of capacity $R(n)$.
    \item For $E$ epochs: sample a mini-batch of size $\beta(n)$ and update $\theta$ by AdamW on Eq.~\ref{eq:loss}.
    \item Step the cosine learning-rate schedule.
  \end{enumerate}
  \end{minipage}}
  \caption{One AlphaZero training iteration for Hex.}
  \label{alg:az}
\end{figure}

\subsection{Board-size scaling rules}
\label{sec:scaling}
The central design choice is that all capacity and budget quantities are functions of $n$.
The per-move simulation budget targets roughly $3.5$ visits per root child, rounded to a
multiple of $50$ and floored at $100$:
\begin{equation}
  N(n) \;=\; \max\!\Big(100,\; 50\cdot\mathrm{round}\!\big(\tfrac{3.5\,n^2}{50}\big)\Big).
  \label{eq:sims}
\end{equation}
At play time the agent uses $2N(n)$ simulations, thinking harder against an opponent than
during data generation. Network capacity steps up for larger boards---$128$ channels and
$5$ residual blocks for $n\le 9$, and $256$ channels and $8$ blocks for $n\ge 11$---because
a deeper trunk is needed for the value head's receptive field to span full-board
connectivity. The number of self-play games per iteration $G(n)$, the replay capacity
$R(n)$ (sized to hold roughly seventy iterations of data), and the mini-batch size
$\beta(n)$ likewise grow with $n$. Table~\ref{tab:scaling} lists the resulting schedule.

\begin{table}[h]
  \caption{Board-size scaling schedule. Simulation budget from Eq.~\eqref{eq:sims};
  inference budget is twice the self-play budget. Network sizes and self-play volume step
  up at $n\ge 11$.}
  \label{tab:scaling}
  \centering
  \begin{tabular}{lcccccc}
    \toprule
    Board $n$ & Self-play sims $N$ & Inference sims & Net (ch$\times$blk) & Games/iter & Params (M) \\
    \midrule
    $7$  & $150$  & $300$  & $128\times5$ & $16$ & --- \\
    $9$  & $300$  & $600$  & $128\times5$ & $16$ & --- \\
    $11$ & $400$  & $800$  & $256\times8$ & $16$ & $9.51$ \\
    $13$ & $600$  & $1200$ & $256\times8$ & $24$ & $9.55$ \\
    $15$ & $800$  & $1600$ & $256\times8$ & $32$ & $9.61$ \\
    $17$ & $1000$ & $2000$ & $256\times8$ & $32$ & $9.69$ \\
    \bottomrule
  \end{tabular}
\end{table}

\paragraph{Design rationale.} The constant $3.5$ in Eq.~\eqref{eq:sims} is the load-bearing
choice. With far fewer visits per child, most children of the root receive zero or one
visit and the visit-count policy $\pi$ degenerates into noise, so the cross-entropy target
in Eq.~\eqref{eq:loss} teaches the network nothing and the bootstrap stalls---the failure
mode we observed when a fixed small budget was reused on large boards. A much larger
constant would make self-play prohibitively slow. Tying capacity, data volume, and search
budget together with a single input also removes the practitioner's most common error:
forgetting to retune one of them when changing $n$.

\subsection{Practical refinements}
Five additions, each standard in strong implementations, complete the system.
\textbf{(1) Symmetry augmentation:} in the perspective frame the rhombus has one
non-trivial symmetry, a $180^\circ$ rotation that maps Red's goal to itself; every
sample is duplicated with its rotated state and policy, doubling data for free.
\textbf{(2) Sub-tree reuse:} after a move, the chosen child's sub-tree is kept as the new
root rather than rebuilt, roughly halving effective search cost in deep positions.
\textbf{(3) Calibrated resignation:} when the root value stays below $-0.95$ for several
consecutive moves the game is conceded to save compute, while a $10\%$ fraction of games
never resign so value targets stay calibrated.
\textbf{(4) Mixed precision:} \texttt{bfloat16} autocast accelerates GPU training and
inference with no loss scaler, since bf16 shares fp32's exponent range.
\textbf{(5) AdamW with cosine annealing} replaces plain Adam plus a manual $L_2$ term,
decoupling weight decay from the adaptive step.

\section{Experiments}
\label{sec:experiments}

\paragraph{Setup.} All training was performed from scratch on a single CUDA GPU using the
C++ MCTS extension. We train independent agents per board size; results below cover
$11\times11$, $13\times13$, and $15\times15$, with longer runs configured for
$15\times15$ and $17\times17$. The opponent for periodic evaluation is a uniform-random
player: across $20$ games (alternating colours) the agent selects moves greedily by visit
count using a reduced budget of $\min(N,15)$ simulations, a deliberately conservative
probe that is cheap enough to run every $20$ iterations. We report the win rate against
random and the two loss components from Eq.~\eqref{eq:loss}. Table~\ref{tab:hparams} lists
the fixed hyperparameters; size-dependent quantities follow Table~\ref{tab:scaling}.

\begin{table}[h]
  \caption{Fixed training and search hyperparameters (size-independent).}
  \label{tab:hparams}
  \centering
  \begin{tabular}{lc@{\hspace{2.5em}}lc}
    \toprule
    Hyperparameter & Value & Hyperparameter & Value \\
    \midrule
    PUCT constant $c_{\text{puct}}$ & $1.5$            & Learning rate (initial) & $10^{-3}$ \\
    Virtual loss                    & $1$              & Weight decay $\lambda$  & $10^{-4}$ \\
    Dirichlet $\alpha$              & $0.3$            & Gradient clip (norm)    & $1.0$ \\
    Dirichlet mix $\varepsilon$     & $0.25$           & Train epochs/iter $E$   & $5$ \\
    Temperature $\tau$              & $1.0$            & Optimizer               & AdamW \\
    Temperature cutoff $T_{\text{cut}}$ & $10$ moves   & LR schedule             & cosine \\
    \bottomrule
  \end{tabular}
\end{table}

\paragraph{Compute.} All experiments were run on a single GPU. Representative throughput:
the $11\times11$ and $13\times13$ agents complete $150$ iterations (about $2{,}400$
self-play games) in a few hours each; the auto-scaled $13\times13$ run extends to $250$
iterations ($\approx 5{,}040$ games) at $600$ simulations per move.
% TODO: exact wall-clock hours per run not logged; fill after instrumented timing.

\section{Results}
\label{sec:results}

\paragraph{Learning against a random opponent.} Table~\ref{tab:main} reports the win rate
versus the random opponent at evaluation checkpoints. On $11\times11$ the agent rises from
chance ($50\%$) at iteration $20$ to $90\%$ by iteration $40$ and reaches $100\%$ by
iteration $80$, fluctuating between $95\%$ and $100\%$ thereafter. The $13\times13$ agent
learns comparably fast, already at $85\%$ by iteration $20$ and saturating near $100\%$.
The $15\times15$ agent, with its larger branching factor, climbs more gradually---from
$55\%$ at iteration $20$ to $90\%$ by iteration $120$---consistent with the expectation
that larger boards need more iterations even with the scaled budget.

\begin{table}[h]
  \caption{Win rate (\%) versus a uniform-random opponent over $20$ games, measured every
  $20$ training iterations with a reduced evaluation budget. Higher is better; $50\%$ is
  chance.}
  \label{tab:main}
  \centering
  \begin{tabular}{lccccccc}
    \toprule
    Board & it.\,20 & it.\,40 & it.\,60 & it.\,80 & it.\,100 & it.\,120 & it.\,140 \\
    \midrule
    $11\times11$ & $50$ & $90$ & $90$ & $\mathbf{100}$ & $95$ & $\mathbf{100}$ & $95$ \\
    $13\times13$ & $85$ & $90$ & $90$ & $\mathbf{100}$ & $95$ & $95$ & $\mathbf{100}$ \\
    $15\times15$ & $55$ & $75$ & $75$ & $70$ & $85$ & $90$ & $80$ \\
    \bottomrule
  \end{tabular}
\end{table}

\paragraph{Extended self-play at higher simulation budgets.} Resuming the $13\times13$
agent under the auto-scaled budget ($600$ simulations per move, $24$ games per iteration)
and training to iteration $240$ holds the win rate at $95$--$100\%$ from iteration $60$
onward (e.g.\ $95\%$ at it.\,60, $100\%$ at it.\,140, $100\%$ at it.\,240), confirming that
the higher-budget regime sustains rather than destabilises the learned policy.

\paragraph{Loss curves.} The value-head mean-squared error settles into the $0.4$--$0.6$
range early and remains stable, indicating calibrated outcome predictions. The policy
cross-entropy decreases slowly---on $13\times13$ from about $4.5$ at iteration $20$ to
$3.9$ at iteration $240$---which is expected: the target $\pi$ is itself a moving, search-
generated distribution over $n^2$ moves, so its entropy floor keeps the absolute loss high
even as the policy sharpens. Table~\ref{tab:loss} summarises representative values.

\begin{table}[h]
  \caption{Representative loss components from Eq.~\eqref{eq:loss} at selected iterations.
  Policy loss is cross-entropy against the MCTS visit distribution; value loss is MSE
  against game outcome.}
  \label{tab:loss}
  \centering
  \begin{tabular}{lcccc}
    \toprule
    Board (run) & Iteration & Policy loss & Value loss & Win\% \\
    \midrule
    $11\times11$ & $20$  & $4.19$ & $0.59$ & $50$ \\
    $11\times11$ & $140$ & $3.98$ & $0.61$ & $95$ \\
    $13\times13$ & $20$  & $4.51$ & $0.51$ & $85$ \\
    $13\times13$ & $140$ & $4.14$ & $0.45$ & $100$ \\
    $13\times13$ ($600$ sims) & $240$ & $3.92$ & $0.49$ & $100$ \\
    $15\times15$ & $20$  & $4.76$ & $0.43$ & $55$ \\
    $15\times15$ & $140$ & $4.23$ & $0.50$ & $80$ \\
    \bottomrule
  \end{tabular}
\end{table}

\paragraph{Effect of the scaling rules.} The qualitative effect of
Section~\ref{sec:scaling} is decisive: with a fixed small simulation budget reused on
$\ge 11\times11$ boards, the visit-count policy target is near-uniform and the agent does
not learn, remaining at chance against random play. Applying Eq.~\eqref{eq:sims} restores
the rapid learning seen in Table~\ref{tab:main}.
% TODO: controlled side-by-side ablation (fixed vs scaled budget, matched compute) not yet
% logged with numbers; report win-rate-vs-iteration for both arms after a paired run.

\section{Discussion}
\label{sec:discussion}

\subsection{Interpretation}
The results support the paper's thesis that, for knowledge-free AlphaZero on Hex, the
governing variable is whether search is deep enough to make the visit-count policy
informative, and that this can be guaranteed by tying the simulation budget to $n^2$. The
ordering of learning speed across boards---fastest on $11\times11$ and $13\times13$, slower
on $15\times15$---matches the prediction that a larger action space dilutes a fixed number
of visits and demands more iterations. That the extended $13\times13$ run remains stable at
a four-fold higher simulation budget indicates the loop is robust to the budget once the
minimum-visit threshold is met, rather than finely balanced on it. The stable, low value
loss alongside a high but slowly falling policy loss is the signature of a healthy
AlphaZero run on a large action space: outcome prediction is easy to calibrate, whereas the
policy chases a non-stationary, high-entropy target.

\subsection{Practical implications}
A single integer suffices to configure the agent for any board in the supported range,
which is precisely what a course or library setting needs: one code path, no per-size
retuning, and a documented failure mode (too-shallow search) with a documented remedy.
The perspective frame and symmetry augmentation show how a small amount of game-structure
awareness---one canonical orientation and one rotation---reduces both the learning problem
and the data cost without importing Hex-specific strategy. The compiled search kernel and
batched, multi-process self-play illustrate that the throughput bottleneck for small-net
AlphaZero is tree-traversal overhead and GPU under-utilisation, both addressable without
changing the algorithm.

\subsection{Relation to prior work}
Our win rates are reported only against a random opponent, whereas specialist Hex engines
such as MoHex are evaluated against each other and against humans \citep{gao2017move,
hayward2019hex}. Defeating random play is a necessary but weak milestone; it confirms the
pipeline learns and is stable, but says little about absolute strength. The clear next
comparison is against the value-based and policy-gradient baselines in the same repository
\citep{mnih2015human,schulman2017proximal} and, ultimately, against a knowledge-rich
solver. We expect the search-based agent to dominate the model-free baselines---this is the
standard finding since ExIt \citep{anthony2017thinking}---and to trail knowledge-rich
engines until Hex-specific inferior-cell and virtual-connection reasoning is added.

\paragraph{Limitations.} Three limitations frame future work. First, evaluation is only
against a random opponent; head-to-head matches against the model-free baselines and a
strong reference engine are needed to quantify absolute strength. Second, the longest runs
for $15\times15$ and $17\times17$ are configured but not yet reported to convergence here,
so conclusions for the largest boards are preliminary. Third, the agent uses no
Hex-specific knowledge, no Gumbel or playout-cap search improvements, and no resignation
ablation; each is a concrete, additive opportunity rather than a structural obstacle.

\section{Conclusion}
\label{sec:conclusion}

We presented a complete, knowledge-free AlphaZero agent for Hex whose defining feature is
that every search and learning quantity---simulation budget, network width and depth,
self-play volume, replay capacity, and batch size---is derived automatically from the board
size through explicit rules. The motivating insight is that learning bootstraps only when
the search visits each root child often enough to make the visit-count policy informative,
a condition that fails silently when a small-board budget is reused on a large board and is
restored by scaling the budget with $n^2$. Around this core we built a reproducible
single-GPU system with a colour-agnostic perspective frame, a compiled C++ MCTS engine with
virtual-loss parallel leaf evaluation, pipelined multi-process self-play, and five standard
refinements. Trained from scratch, the agent reaches $90$--$100\%$ win rates against a
random opponent on $11\times11$ and $13\times13$ within tens of iterations and continues to
improve on $15\times15$.

Three concrete directions follow. First, replace the random-opponent metric with
head-to-head Elo against the repository's DQN, PPO, and REINFORCE baselines and against a
reference Hex engine, to measure absolute rather than relative strength. Second, run the
$15\times15$ and $17\times17$ agents to convergence and conduct a paired, compute-matched
ablation of the scaling rules and each of the five refinements. Third, add modern search
and knowledge improvements---Gumbel-AlphaZero action selection, playout caps, and
Hex-specific inferior-cell pruning---to close the gap to specialist engines. More broadly,
the size-scaling principle suggests a path toward agents that self-configure across a
family of related tasks from a single structural parameter, reducing the manual tuning that
makes self-play reinforcement learning fragile.

\begin{ack}
This work was developed as part of the reinforcement learning coursework at the University
of Applied Sciences Technikum Wien (FHTW). The author declares no competing interests and
received no external funding.
\end{ack}

\section*{References}
{\small
\begin{thebibliography}{99}

\bibitem[Nash(1952)]{nash1952hex}
  Nash, J.\ (1952).
  Some games and machines for playing them.
  {\it RAND Corporation Technical Report}, D-1164.

\bibitem[Silver et al.(2016)]{silver2016mastering}
  Silver, D., Huang, A., Maddison, C.J., Guez, A., Sifre, L., van den Driessche, G.,
  et al.\ (2016).
  Mastering the game of Go with deep neural networks and tree search.
  {\it Nature}, 529(7587), 484--489.

\bibitem[Silver et al.(2017)]{silver2017mastering}
  Silver, D., Schrittwieser, J., Simonyan, K., Antonoglou, I., Huang, A., Guez, A.,
  et al.\ (2017).
  Mastering the game of Go without human knowledge.
  {\it Nature}, 550(7676), 354--359.

\bibitem[Silver et al.(2018)]{silver2018general}
  Silver, D., Hubert, T., Schrittwieser, J., Antonoglou, I., Lai, M., Guez, A., et al.\ (2018).
  A general reinforcement learning algorithm that masters chess, shogi, and Go through
  self-play.
  {\it Science}, 362(6419), 1140--1144.

\bibitem[Schrittwieser et al.(2020)]{schrittwieser2020mastering}
  Schrittwieser, J., Antonoglou, I., Hubert, T., Simonyan, K., Sifre, L., Schmitt, S.,
  et al.\ (2020).
  Mastering Atari, Go, chess and shogi by planning with a learned model.
  {\it Nature}, 588(7839), 604--609.

\bibitem[Anthony et al.(2017)]{anthony2017thinking}
  Anthony, T., Tian, Z., \& Barber, D.\ (2017).
  Thinking fast and slow with deep learning and tree search.
  {\it Advances in Neural Information Processing Systems}, 30.

\bibitem[Gao et al.(2017)]{gao2017move}
  Gao, C., Hayward, R., \& M\"uller, M.\ (2017).
  Move prediction using deep convolutional neural networks in Hex.
  {\it IEEE Transactions on Games}, 10(4), 336--343.

\bibitem[Hayward \& Toft(2019)]{hayward2019hex}
  Hayward, R.B., \& Toft, B.\ (2019).
  {\it Hex: The Full Story}.
  CRC Press.

\bibitem[Kocsis \& Szepesv\'ari(2006)]{kocsis2006bandit}
  Kocsis, L., \& Szepesv\'ari, C.\ (2006).
  Bandit based Monte-Carlo planning.
  {\it European Conference on Machine Learning (ECML)}, 282--293.

\bibitem[Browne et al.(2012)]{browne2012survey}
  Browne, C., Powley, E., Whitehouse, D., Lucas, S., Cowling, P.I., Rohlfshagen, P.,
  et al.\ (2012).
  A survey of Monte Carlo tree search methods.
  {\it IEEE Transactions on Computational Intelligence and AI in Games}, 4(1), 1--43.

\bibitem[Chaslot et al.(2008)]{chaslot2008parallel}
  Chaslot, G.M.J.-B., Winands, M.H.M., \& van den Herik, H.J.\ (2008).
  Parallel Monte-Carlo tree search.
  {\it International Conference on Computers and Games}, 60--71.

\bibitem[Tian et al.(2019)]{tian2019elf}
  Tian, Y., Ma, J., Gong, Q., Sengupta, S., Chen, Z., Pinkerton, J., \& Zitnick, C.L.\ (2019).
  ELF OpenGo: An analysis and open reimplementation of AlphaZero.
  {\it International Conference on Machine Learning (ICML)}, 6244--6253.

\bibitem[Wu(2019)]{wu2019accelerating}
  Wu, D.J.\ (2019).
  Accelerating self-play learning in Go.
  {\it arXiv preprint} arXiv:1902.10565.

\bibitem[He et al.(2016)]{he2016deep}
  He, K., Zhang, X., Ren, S., \& Sun, J.\ (2016).
  Deep residual learning for image recognition.
  {\it IEEE Conference on Computer Vision and Pattern Recognition (CVPR)}, 770--778.

\bibitem[Ioffe \& Szegedy(2015)]{ioffe2015batch}
  Ioffe, S., \& Szegedy, C.\ (2015).
  Batch normalization: Accelerating deep network training by reducing internal covariate
  shift.
  {\it International Conference on Machine Learning (ICML)}, 448--456.

\bibitem[Loshchilov \& Hutter(2019)]{loshchilov2019decoupled}
  Loshchilov, I., \& Hutter, F.\ (2019).
  Decoupled weight decay regularization.
  {\it International Conference on Learning Representations (ICLR)}.

\bibitem[Mnih et al.(2015)]{mnih2015human}
  Mnih, V., Kavukcuoglu, K., Silver, D., Rusu, A.A., Veness, J., Bellemare, M.G.,
  et al.\ (2015).
  Human-level control through deep reinforcement learning.
  {\it Nature}, 518(7540), 529--533.

\bibitem[Schulman et al.(2017)]{schulman2017proximal}
  Schulman, J., Wolski, F., Dhariwal, P., Radford, A., \& Klimov, O.\ (2017).
  Proximal policy optimization algorithms.
  {\it arXiv preprint} arXiv:1707.06347.

\bibitem[Williams(1992)]{williams1992simple}
  Williams, R.J.\ (1992).
  Simple statistical gradient-following algorithms for connectionist reinforcement
  learning.
  {\it Machine Learning}, 8(3), 229--256.

\bibitem[Watkins \& Dayan(1992)]{watkins1992qlearning}
  Watkins, C.J.C.H., \& Dayan, P.\ (1992).
  Q-learning.
  {\it Machine Learning}, 8(3), 279--292.

\end{thebibliography}
}

\appendix

\section{Implementation notes}
\label{app:impl}
The agent interface is a single callable \texttt{agent(board, action\_set)} returning a
$(\text{row},\text{col})$ move. On the first move of an empty board the agent plays a random
interior cell, avoiding the outer ring whose corner/edge openings are weak in Hex; all
subsequent moves use the full inference-time MCTS budget $2N(n)$. Models are persisted per
board size with their channel and block counts so they reload with the correct architecture.
Self-play workers run network inference on CPU to avoid contending with the GPU training
process, communicating weights and samples through multiprocessing queues with the
file-system tensor-sharing strategy to avoid file-descriptor exhaustion. The repository also
contains model-free baselines---tabular Q-learning, DQN (with an optional minimax variant),
REINFORCE, and PPO---behind the same interface for comparison.

\section{Reproducibility}
\label{app:repro}
Training is launched via \texttt{python train.py --method alphazero --size N
--iterations K}, with all size-dependent quantities defaulting to the rules of
Section~\ref{sec:scaling} unless overridden. A guard warns when the requested parallel
leaf count exceeds $N/7$, which would make the search too shallow to learn. Checkpoints are
written every $20$ iterations, enabling resumption at a higher simulation budget as used for
the extended $13\times13$ run.

\end{document}
